\chapter{Background \& Related Work}

\section{Graph Partitioning}
The general problem of exact graph partitioning is known to be NP-complete \cite{garey1974some}, as a result of this severe restriction, the large body of research mostly focuses on approximate solutions to this problem over exact solutions, perhaps the best known one being the Kernighan-Lin algorithm \cite{kernighan1970efficient}. Most graph clustering algorithms are based upon a few key concepts such as heuristics \cite{kernighan1970efficient}, agglomerative clustering \cite{newman2004fastcommunity, clauset2004finding} and divisive clustering \cite{newman2004findingeval} although other techniques such as spectral clustering, k-means clustering, do exist \cite{hopcroft2003naturalcommunity, alzate2008multiway, lang2004flow}. The full literature on graph clustering algorithms is extensive and cannot be expressed concisely, however, a large amount of these algorithms have been studied and summarized in various surveys and studies \cite{schaeffer2007graph, abbas2018streaming, leskovec2010empirical}. The specific area concerning real-time graph clustering algorithms mostly consists of simple, computationally inexpensive, heuristic based methods \cite{stanton2012streaming} as mentioned previously, although some exotic methods have also been presented \cite{zanghi2008fast}. 

\section{Incremental Computation}
Incremental Computation has been an niche, unique field with relatively few published literature, most published research is based upon Umut Acar's PhD thesis on "Self Adjusting Computations" \cite{acar2005self}, later developments such as Adapton \cite{hammer2014adapton} build upon this work. The idea behind this class of incremental computation is somewhat simple. Adapton for example performs general purpose incremental computations through building a computation graph, whenever data in this computation graph is changed, the child nodes in this computation graph are ``dirtied'', whenever the result of this computation graph is requested by the application, these ``dirty'' nodes are recomputed. The fundamental difference between Differential Dataflow and Adapton is that Adapton is a general purpose incremental computation framework while Differential Dataflow is optimised for writing dataflow programs and the open source implementation of Differential Dataflow is built on a distributed low latency, cyclic dataflow programming layer called `Timely Dataflow' \cite{murray2013naiad}. 


\section{Incremental View Maintenance}
Incremental View Maintenance algorithms are of concern when discussing Differential Dataflow as they both aim to perform similar objectives. However, it should be noted that existing IVM techniques either perform too much work, maintain too much state or limit expressiveness. Take for consideration the classic incremental view maintainence algorithm by Gupta \cite{gupta1993maintaining}, this algorithm has some execution such that it over-estimates the number of invalidated tuples and performs a large amount of work to correct this error. Consider Nigam's work on the extended PSN algorithm \cite{nigam2012maintaining}, this implementation stores tuples along with the full set of tuples used to derive it, this is obviously restrictive on large networks. If we take into account later developments such as DBToaster \cite{ahmad2012dbtoaster} that has spearheaded IVM algorithm research since 2009, it still suffers from the inability to perform IVM on iterative workloads. 

\section{Incremental Dataflow}
While there have been several papers and software solutions to the problem of incrementally processing dataflow programs, they have been mostly simple modifications to existing dataflow systems. Incoop \cite{bhatotia2011incoop}, builds upon Acar's previous work on "Self Adjusting Computations" \cite{acar2005self} to bring incremental processing to the MapReduce paradigm, on closer inspection of Incoop, one would notice that the suggested system is not much more than a simple caching and reuse of values. There are plenty of obvious drawbacks when using such a system, such as the inability to process iterative algorithms and the inability to maintain low latency. 

\section{Iterative Dataflow}
Various attempts have been made to enhance the MapReduce \cite{dean2004mapreduce} paradigm to support iterative workloads. The core principle here is to repeatedly apply MapReduce until a defined stopping criteria is met. Apart from the fact that this modification can not deal with incremental workloads, the implementations of iterative MapReduce introduce additional problems regarding data consistency after node failures. Consider for example the instance of a node failure in a MapReduce system where an iterative workload is in the process of being computed, data consistency issues may arise if the nodes execute on data at different iteration points after the restoration of the failed node, these issues can be addressed but introduce additional system complexity and latency. 


% Show differential dataflow here?

\section{Modularity Optimisation}
Modularity optimisation was initially proposed by Newman \cite{newman2004findingeval}, while algorithms like Louvain \cite{blondel2008fast} exist to perform this task, it is important to note that exact modularity optimisation, like graph partitioning is NP-hard. It is because of this fundamental limitation that the algorithms that have been researched and presented in literature are only heuristic based methods that do not aim for optimal modularity optimisation but rather to simply get close to optimal results. Unfortunately, this reliance on a global metric for the computation of modularity is why parallel and dataflow implementations of modularity optimisations are hard to implement and why they suffer from convergence issues, deadlocks and negative gain scenarios. 

\section{The Louvain Algorithm}
Louvain \cite{blondel2008fast} is an agglomerative clustering algorithm based upon optimising modularity, the only objective of this algorithm is to optimise $Q=\frac{1}{2 m} \sum_{i, j}\left[A_{i j}-\frac{k_{i} k_{j}}{2 m}\right] \delta\left(c_{i}, c_{j}\right)$. The algorithm is initiated through allocating each node into its own unique community, then two phases are performed in iteration until a fixed point is reached or a defined modularity threshold is not met. The first phase is to simply iterate over the nodes and calculate the gain in modularity by moving the node into the communities of its neighbours, this change in modularity can be calculated through the equation: $$\Delta Q=\left[\frac{\sum_{i n}+k_{i, i n}}{2 m}-\left(\frac{\sum_{t o t}+k_{i}}{2 m}\right)^{2}\right]-\left[\frac{\sum_{i n}}{2 m}-\left(\frac{\sum_{t o t}}{2 m}\right)^{2}-\left(\frac{k_{i}}{2 m}\right)^{2}\right]$$. After the completion of this evaluation for a given vertex $i$, a neighbour $j$ is chosen such that it represents the maximum modularity gain by moving $i$ into the communities of one of its neighbours, in the case where the modularity gain is zero, no community transfers take place. This step is repeated until a fixed point or a defined stopping criteria is met. It is important to note that Louvain does unfortunately evaluate modularity needlessly for certain nodes, contributing to a performance degradation. The second phase of this algorithm is to simply aggregate all nodes within the same community to a new node, as mentioned previously these two phases are repeated until a stopping criteria is met. It is important to note that Louvain \cite{blondel2008fast} is able to offer significantly faster clustering with a runtime complexity of $\mathcal{O}(n*log(n))$, this proved to be a clear advantage at the time where most algorithms were $\mathcal{O}(n^p)$ at least, where $p$ was some integer greater than 1. The algorithm however is not without faults, unfortunately one disadvantage of using the above mentioned aggregation step is that it is greedy and only has a limited view into how nodes will be clustered globally and throughout the various iterations, as a result of this it may not yield optimal modularity for the clustered graph. To demonstrate one failing of Louvain, consider the case where it may be beneficial to move a node from a cluster to a newly formed cluster, unfortunately due to the aggregation step, this can never happen. Unfortunately, Louvain has no known upper bound for the number of iterations that have to be performed in order to obtain convergence, this is why simpler methods such as meeting a modularity threshold or even simply capping the number of iterations to some fixed integer is typically used. 

\section{Limits of modularity}
It is important to note that optimisation of modularity is not a guarantee of the detection of community structures, this is further explored in \cite{fortunato2010community}. It has been shown that there even exists large values of measures of modularity in completely random graphs \cite{guimera2004modularity, reichardt2006statistical}, this obviously presents a major limitation in evaluating the true community structure through measuring modularity alone. Fortunato notes that this is as a result of random fluctuations in the distribution of edges, he states that these fluctuations determine concentrations of links in subsets of the graph which then falsely may appear to be community structures. As a result of these findings regarding modularity measurement, we cannot definitively state if a measurement of modularity corresponds to some community structure, however Fortunato has determined that a z-score derived from random rewiring of the graph could be used to obtain some probabilistic measure of community detection. To make matters worse, Fortunato has noted a fundemental problem at the core of modularity optimisation techniques and that is the fact that modularity optimisation has a resolution limit. To be more precise, if the maximum modularity produces communities of size $\sqrt{m}$ or smaller, one cannot know if that particular community is a single community or in fact made up of smaller weakly linked communities. While there have been some attempts to address this issue, our implementation exclusively looked at the modularity definition present in the Louvain method \cite{blondel2008fast}, despite its inherent problems simply as a result of that definition of modularity being more commonly used. 
